\documentclass[11pt]{article}
\usepackage[margin=2.3cm]{geometry}
\usepackage{amsmath,amssymb,booktabs}
\usepackage[hidelinks]{hyperref}
\usepackage{microtype}
\usepackage{url}
\DeclareUrlCommand\lean{\urlstyle{tt}\def\UrlBreaks{\do\.\do\_\do\/}}
% the claim lists of a statement and their byte encodings
\newcommand{\Cx}{C_{\mathrm{xmss}}}
\newcommand{\Cs}{C_{\mathrm{sphincs}}}
\newcommand{\enc}{\mathrm{enc}}
% the counts of a node's work (Section~\ref{sec:guest}, the budget)
\newcommand{\cnt}[1]{n_{\mathrm{#1}}}
\emergencystretch=3em
\setlength{\parskip}{0.35em}
\title{What is proved about leanVM}
\author{leanerVM, branch \texttt{mumbai}, about leanVM \texttt{c445e3d5}}
\date{29 October 2026}
\begin{document}
\maketitle

\noindent This document states, in plain mathematics, every end statement the repository proves:
leanVM is knowledge sound and complete; its leanISA guest program, which aggregates signatures,
data-availability claims and other proofs, is sound and complete as an aggregation step; leanXMSS
and leanSPHINCS are strongly unforgeable; and LeanDA's commitments are position binding, and code
binding for arguments whose relation is evaluated over the model's permutation. Each is a Lean~4
proposition proved using only Lean's three axioms (\lean{propext}, \lean{Classical.choice},
\lean{Quot.sound}): no \lean{sorry}, no \lean{native_decide}, no added axiom. The system is
leanVM at commit \texttt{c445e3d5} of its \texttt{main} branch.

\paragraph{Where to read it.} The statements live in one folder, \lean{Statement/}: fifteen files
of definitions and end statements ($7421$ lines with their documentation), which import Mathlib,
VCVio and CompPoly and nothing of the proofs. Beside them sit the honest prover the completeness
statements name (\lean{Statement/Prover/}, $2477$ lines; it is proved to make the specification's
honest moves, so a reviewer only needs to see that it is a genuine prover), the guest program's
generated bytecode (\lean{Statement/Data/}), and proved \lean{@[csimp]} lemmas for speed only
(\lean{Statement/Fast/}). \textbf{The entry point is \lean{Statement/Theorems.lean}}: the
structure \texttt{Statement.Theorems decode} has one field per end statement, and
\lean{lean_vm/LeanVM.lean} proves them all at once,
\texttt{LeanVM.theorems\_hold : Statement.Theorems LeanVM.Guest.End.decode}, its axioms pinned by
the build.
\begin{center}\small
\begin{tabular}{lll}
\toprule
Field & End statement & Section\\\midrule
\lean{knowledgeSoundness} & \lean{LeanVM.KnowledgeSoundness} & \ref{sec:zkvm}\\
\lean{completeness} & \lean{LeanVM.Completeness} & \ref{sec:zkvm}\\
\lean{leanXmss} & \lean{LeanXmss.Rpm.LeanXmssStronglyUnforgeable} & \ref{sec:schemes}\\
\lean{leanSphincs} & \lean{Sphincs.Lifetimes.LeanSphincsStronglyUnforgeable} & \ref{sec:schemes}\\
\lean{leanDaComplete} & \lean{LeanDa.PerfectlyCompleteRpm} & \ref{sec:schemes}\\
\lean{leanDaPositionBinding} & \lean{LeanDa.PositionBindingRpm} & \ref{sec:schemes}\\
\lean{leanDaCodeBinding} & \lean{LeanDa.CodeBindingRpm} & \ref{sec:schemes}\\
\lean{guestKnowledge} & \texttt{LeanVM.Node.GuestKnowledge guest decode} & \ref{sec:guest}\\
\lean{digestBinds} & \lean{LeanVM.Node.DigestBinds} & \ref{sec:guest}\\
\lean{batchSound} & \lean{LeanVM.Guest.Batch.BatchingSound} & \ref{sec:guest}\\
\lean{guestCompleteCovered} & \lean{LeanVM.Guest.Budget.GuestCompleteCovered} & \ref{sec:guest}\\
\lean{guestChildComplete} & \lean{LeanVM.Guest.Budget.GuestChildComplete} & \ref{sec:guest}\\
\bottomrule
\end{tabular}
\end{center}
The twelfth file, \lean{Statement/Corollaries.lean}, is optional reading: it states three
consequences of the fields, \lean{AggregateSound}, \lean{GuestCompleteRecursion} and
\lean{GuestCompleteLeaf} (Section~\ref{sec:guest}), proved from the fields alone by
\lean{LeanVM.corollaries_of_theorems}. A few theorems cited below are proved in the proof
projects, but their statements are not part of \lean{Statement/}: \lean{decode_binds}, \lean{guest_sound},
\lean{aggregate_knowledge}, \lean{guest_complete_numeric}, \lean{aggregate_complete_of},
\lean{completenessRpm} and \lean{leanDa_codeBinding_core}.

\section{The setting}

\paragraph{Hashing.} leanVM hashes with Keccak-256, the sponge over the permutation
Keccak-f[1600] (rate $136$ bytes, capacity $512$ bits, padding \texttt{0x01}\,\dots\,\texttt{0x80},
zero initial state, the first four lanes as digest). The statement writes a state as $25$ lanes of
$64$ bits throughout and defines Keccak-256 once, over any permutation (\lean{keccak256With});
the deployed Keccak-256 is this sponge over Keccak-f[1600]. leanVM's Fiat--Shamir transcript is a
duplex that calls Keccak-f[1600] directly, overwriting the first $16$ lanes of the state.

\paragraph{The random permutation model (RPM).} Keccak-f[1600] is modelled as a uniformly random
permutation $\pi$ of $\{0,1\}^{1600}$, sampled lazily (\lean{Common.lazyRandomPermutation},
\lean{Statement/RandomPermutation.lean}). The model's oracles are $\pi$ and $\pi^{-1}$ and
nothing else (\lean{Common.RandomPermutationOracles}). No hash is an oracle or a parameter of
the statement: every algorithm (verifier, prover, scheme, argument) is written over one
permutation and computes Keccak-256 as the sponge over it (\lean{keccak256With}), and its duplex
calls it; the model gives every party $\pi$ (\lean{Common.queryPermutation}), the deployed
instance Keccak-f[1600]. Each experiment is one game run against the lazily sampled permutation
(\lean{Common.runWithRandomPermutation}), the adversary's queries recorded where an extractor reads
them (\lean{Common.loggedOracles}). \emph{The assumption is that the
deployed Keccak-f[1600], a fixed permutation, behaves like a random one}, even to programs that
compute it: the \texttt{KECCAK} instruction a proved program executes computes the deployed
Keccak-f[1600], so in the model the transcript's permutation and a program's are independent,
while deployed they are the same function. The model thus excludes, by construction, programs that
evaluate their own proof's Fiat--Shamir duplex, the setting of the known attacks on Fiat--Shamir
for GKR-based systems (Khovratovich, Rothblum and Soukhanov, 2025). The adversary is classical,
and its code is fixed before $\pi$ is sampled (no advice that depends on the permutation). Every end statement with
an adversary is in this model; \lean{Statement/} contains no random oracle. Queries are counted in
one of two ways:
\begin{itemize}
\item for leanVM and its guest program (Sections~\ref{sec:zkvm} and~\ref{sec:guest}), $Q$ bounds
the adversary's own queries to $\pi$ and $\pi^{-1}$ on every path; the verifier's calls are bounded
in the proofs and enter the constants;
\item for the signature schemes and LeanDA (Section~\ref{sec:schemes}), $Q$ bounds every
permutation call of the game on every path, the scheme's own included; a Keccak-256 evaluation of
$n$ bytes costs $\lfloor n/136\rfloor + 1$ calls.
\end{itemize}
Adversaries may flip private coins, which are neither counted nor recorded.

\paragraph{How the proofs reach the RPM.} Each proof is first carried out with independent random
oracles (Keccak-256 a random function on byte strings and, for leanVM and LeanDA, Keccak-f a random
function on states), then moved to the RPM by the project \lean{common}: the sponge over $\pi$ is
indifferentiable from a random oracle, so an experiment of $q$ permutation calls changes
probability by at most $q^2/2^{509}$, or $q^2/2^{508}$ with the duplex's direct calls
(\lean{Common.keccak256_transfer}, \lean{keccak256_transfer2}). The proofs reason with Keccak-256 as a call or parameter of its own, over
the same $\pi$; with that hash Keccak-256 over the permutation these forms are the statement's
algorithms and games, with the same distribution and the same permutation calls (the
\lean{WithHashBridge} files of each project, \lean{Common.PiOnly}). These are proof steps; their
costs are included in the bounds below.

\section{The zkVM: knowledge soundness and completeness}\label{sec:zkvm}

\paragraph{What a proof claims.} A proof claims that a program $P$ (leanVM bytecode) runs to
completion on a public input $x$. A \emph{valid execution} of $P$ on $x$ (\lean{ValidExecution},
\lean{Statement/LeanIsa.lean}) is a memory of $2^\kappa$ words, $2^{16} \le 2^\kappa \le 2^{32}$,
holding $x$ in words $0$ and $1$, together with a number of steps that run $P$ on it from
$(\mathit{pc}, \mathit{fp}) = (1, 1)$ to the final state $(g^{N-1}, 1)$, where $N$ is the
bytecode's size: the last bytecode slot, in \texttt{main}'s frame. The machine halts when both
registers equal the final state's, checked before each step. This is exactly the state the
verifier's state bus checks for, so knowledge soundness needs no condition on the program.

\paragraph{The extractor.} Knowledge soundness is stated for \emph{one fixed} extractor,
\lean{transcriptExtractor} (\lean{Statement/LeanVm.lean}). It is straight-line: it reads the proof and
the adversary's transcript and nothing else; it neither queries $\pi$, nor rewinds the adversary,
nor sees the verifier's queries. From the transcript's permutation queries (its coins left out,
\lean{permutationQueries}) it collects the Keccak-256 evaluations whose
whole padded sponge path appears as a chain of forward queries from the zero state
(\lean{keccakEvaluationsOfTranscript}). It reads the announced sizes and the stack's Merkle root off the proof,
inverts the Merkle tree through those evaluations to recover the committed word, lists the stacks
whose Reed--Solomon encoding agrees with the word on at least the Johnson threshold of columns (at
most $2^{18}$ of them, \lean{LeanVM.candidates_ncard_le}), and returns a candidate's memory
image whose halting run is a valid execution, if one exists.

It is an algorithm: every definition it uses is computable and deterministic. It lists the
candidates lane by lane with the Guruswami--Sudan list decoder over $\mathbb{F}_{2^{64}}$
(CompPoly's Lee--O'Sullivan interpolation and Roth--Ruckenstein root finding, with Shoup's trace
splitting for the field roots), and it is proved to list every candidate. It runs each
candidate for at most as many steps as the announced tables have rows (\lean{stepCap}), which a
valid execution never exceeds (proved). That it runs in polynomial time is read off its loops
(a cost table in \lean{Statement/Extractor.lean}), not proved; on the words with errors an
adversary chooses it is infeasible at deployed sizes. It gives straight-line knowledge, not a
concrete-time extraction: the collision clauses below carry no concrete reduction cost.

\paragraph{Knowledge soundness} (\lean{LeanVM.KnowledgeSoundness}). In the experiment
(\lean{soundnessExperiment}), an adversary $\mathcal{A}$ with oracle access to $\pi$ and $\pi^{-1}$
outputs a program $P$, a public input $x$ and a proof, and the verifier checks it with Keccak-256
the sponge over $\pi$ and Keccak-f the permutation $\pi$. For every adversary making at most
$Q \le 2^{128}$ queries to $\pi$ and $\pi^{-1}$ on every path,
\[
\boxed{\ \Pr\left[\begin{array}{l}\text{the verifier accepts } (P,x,\text{proof}) \text{ and the}\\
\text{extractor outputs no valid execution of } P \text{ on } x\end{array}\right]
\le (Q+2^{14})\cdot 2^{-128} + 5\,(Q+2^{14})^2\cdot 2^{-256}\ }
\]
(\lean{LeanVM.knowledgeSoundness_holds}). The constant $2^{14}$ absorbs the verifier's $12006$
calls and two further terms of at most $2^{-128}$ each: the transfer to the RPM, and the
verifier's own $12006/2^{256}$.

\paragraph{Completeness} (\lean{LeanVM.Completeness}). The statement names its honest prover,
\lean{LeanVM.prover}: the executable \lean{LeanVM.Prover.prove} (\lean{Statement/Prover/}), its
Keccak-256 and Keccak-f[1600] calls answered here from $\pi$; with the deployed Keccak it is
\texttt{lake exe leanvm-prove}. It runs the execution once to obtain the witness, then plays the
honest moves against the verifier's own program. Its only search is the proof-of-work nonce
search within its fuel, and it sees hash answers only as the challenges the verifier samples. On
the test corpus ($17$ programs at $4$ rates) its proofs are byte-identical to the Rust prover's,
and both verifiers reject every mutation of them.

For every program $P$, public input $x$ and valid execution of $P$ on $x$ of committed size at
most $2^{28}$ cells, at every inverse rate $2^\rho$ with $1 \le \rho \le 4$, trying $2^{25}$ nonces
per proof of work,
\[
\boxed{\ \Pr[\text{the verifier rejects the honest proof}] \le 2^{-256}\ }
\]
(\lean{LeanVM.completeness_holds}, through \lean{prove_eq_honestComp}: the prover is the
specification's honest prover at the execution's witness). Completeness is not perfect only
because the proof of work is bounded: for any fuel $f \le 2^{64}$ the rejection probability is at
most
\[
\delta(f) = Q_f\,(1-2^{-17})^{f} + (Q_f\cdot 2^{32})^2/2^{508}, \qquad
Q_f = 12008 + 48017\,(2^{27} + 2f + 2)
\]
(\lean{LeanVM.Rpm.completenessRpm}). The \emph{committed size} (\lean{Execution.committedSize}) of an
execution is
$4\cdot 2^{\kappa} + 2^{\ell} + 15h_{\mathrm{xor}} + 15h_{\mathrm{mul}} + 8h_{\mathrm{set}} +
15h_{\mathrm{deref}} + 14h_{\mathrm{jump}} + 26k + 2^6 k$, where $2^\kappa$ is the memory size,
$2^\ell$ the bytecode size, $h_{op}$ the least power of two at least the number of executed
instructions of opcode $op$, and $k$ the same for \texttt{KECCAK}, at least $2^9$.

\section{The signature schemes and LeanDA}\label{sec:schemes}

Here $Q$ counts every permutation call of the game, and every bound holds for every $Q$ (a bound
above one says nothing).

\subsection*{The signature schemes}

\paragraph{leanXMSS} (\lean{LeanXmss.Rpm.LeanXmssStronglyUnforgeable},
\lean{Statement/LeanXmss.lean}): a single key is strongly unforgeable under chosen-message attack.
The game samples the key's seed, generates the keys, and runs the adversary with $\pi$,
$\pi^{-1}$ and a signing oracle; the adversary picks each request's leaf index and wins
only if no index is requested twice, and its forgery wins if it verifies and is not a pair the
oracle returned. For every $Q \ge 1$, the forgery probability is at most
\[
\boxed{\ (Q + 1)\cdot 2^{-127}\ }
\]
(\lean{XmssSecurity.leanXmss_rpm_security}). $Q$ includes key generation, every signing request
and the final verification at their worst case: a signature grinds through up to $2^{32}$
randomizers in each of its two layers, so one request alone may cost over $2^{33}$ calls.

\paragraph{leanSPHINCS} has the same game and the same bound, $(Q+1)\cdot 2^{-127}$ for every
$Q \ge 1$ (\lean{Sphincs.Lifetimes.LeanSphincsStronglyUnforgeable}, \lean{Statement/LeanSphincs.lean}), with the deterministic signer and at most the lifetime's number of signing
requests, at each of the seven proved lifetimes (subtree heights $26, 20, 14, 13, 12, 10, 8$;
\lean{Sphincs.Lifetimes.requestedSecurityRpm}). A signature costs about $2^{35}$ calls in the worst
case, so with every request of the largest lifetime used the bound is about $2^{-62}$. Both
schemes are stated through one game, \lean{Common.SufCma.StronglyUnforgeable}, and share one tweakable hash
(\texttt{tweak.rs}), hash chain, path fold and tree.

\subsection*{LeanDA}

LeanDA commits to a matrix of blobs so that clients can check random pieces of it. Its statements are
Benedikt Wagner's framework, theorems and proofs, ported and instantiated with LeanDA
(\lean{Statement/LeanDa.lean}). Everyone in the games (the adversary, the argument, the verifier)
queries $\pi$ and $\pi^{-1}$ alone and hashes with Keccak-256 over $\pi$.

\paragraph{The matrix.} The matrix has $r$ rows, $1 \le r \le 1024$, one blob each. A blob is a
payload of $2^{14}$ field elements, Reed--Solomon encoded to $2^{15}$ (\texttt{code r}). Each
encoded row is cut into $128$ \emph{cells}, so the matrix has $128\,r$ positions.

\paragraph{The commitment.} The committer hashes the cells into two Merkle trees, one over rows
and one over columns, and publishes $\mathit{root} = H(\mathit{rowRoot} \,\|\, \mathit{colRoot})$
(\lean{commitCells}). To \emph{open} a position, it gives the cell, the row root, and the cell's
path in the column tree; the opening verifies if they hash back to the root (\lean{verifyCellM}).

\paragraph{The proof that the rows are codewords.} A Merkle root alone does not show that the rows
are codewords. So the committer also gives a proof, from some \emph{argument} (\lean{OracleArgument}: a
prover and a verifier that may query the oracles), that the cells satisfy LeanDA's relation
(\lean{relationM}):
\begin{itemize}
\item the cells hash to the root, and
\item every row passes the \emph{membership test}: from the root, the duplex derives $14$
challenges, which define a codeword $L$ of the dual code; every row must be orthogonal to $L$. A
row that is not a codeword passes with probability at most $14/2^{192}$.
\end{itemize}

\paragraph{Completeness} (\lean{LeanDa.PerfectlyCompleteRpm}). If the argument is perfectly
complete, every position of an honest commitment opens and verifies, with probability $1$, for
every permutation $\pi$ (\lean{LeanDa.leanDa_perfectlyComplete_rpm}).

\paragraph{Position binding} (\lean{LeanDa.PositionBindingRpm}). No adversary can open one
position to two different cells. For every argument, with $Q$ permutation calls, it succeeds with
probability at most
\[
\boxed{\ (Q+1)^2\cdot 2^{-256}\ }
\]
(\lean{LeanDa.leanDa_positionBinding_rpm}).

\paragraph{Code binding} (\lean{LeanDa.CodeBindingRpm}). Position binding is not enough for
sampling. A committer could commit to rows that are not codewords; each opening would still
verify, but the cells clients receive would not come from one matrix of codewords, whose blobs could then not
be reconstructed. Code binding says this cannot happen.

\emph{The game} (\lean{codeBindingGame}). The adversary outputs a root, a proof, and some opened
positions (each with a cell and an opening). It wins if
\begin{itemize}
\item the argument's verifier accepts the proof,
\item every opening verifies, and
\item no codeword agrees with all the opened cells (\lean{ErasureCode.Consistent}).
\end{itemize}

\emph{The assumption.} Code binding can only be as good as the argument, so the theorem assumes the
argument is \emph{knowledge sound}. That means there is an extractor that, reading only the proof
and the adversary's queries (\lean{StraightLineExtractor}; it does not rewind), finds cells satisfying the
relation whenever the verifier accepts. Precisely: in the game \lean{knowledgeSoundnessGame}, an adversary making
$q \le 2^{128}$ calls makes the verifier accept while the extracted cells fail the relation with
probability at most $(q + c)\cdot 2^{-k}$, for some constants $c$ and $k$.

\emph{The theorem.} Under this assumption, an adversary making $Q$ calls (any $Q$, verifications
included) wins the code-binding game with probability at most
\[
\boxed{\ (Q + c + 2^{22})\cdot 2^{-k} + (Q + 2^{22})\cdot 2^{-188} + (Q + 2^{22})^2\cdot 2^{-247}\ }
\]
(\lean{LeanDa.leanDa_codeBinding_rpm_bound}). At $Q = 2^{64}$ the last two terms are about
$2^{-124}$ and $2^{-119}$.

\emph{Why.} Run the extractor on the adversary; it returns cells $c^*$. One of three things holds.
\begin{enumerate}
\item $c^*$ fails the relation. Then the adversary breaks the argument's knowledge soundness:
the first term.
\item $c^*$ satisfies the relation but has a row that is not a codeword. Then it passed the
membership test by luck, $14/2^{192}$ per root: the second term.
\item $c^*$ is a codeword. Every opened cell hashes to the same root as $c^*$, so unless Keccak-256
has a collision, each opened cell is $c^*$'s cell, and the codeword $c^*$ agrees with all of them:
the adversary loses. Collisions and other bad hash events give the third term.
\end{enumerate}
This proof needs no rewinding. Wagner's original proof, with a forking lemma, is also formalized
(\lean{LeanDa.leanDa_codeBinding_core}, for arguments that do not query the oracles); it gives
the weaker $(Q + c + 2^{22})\cdot 2^{-k} + (Q + 2^{20})\cdot 2^{-93} + (Q + 2^{22})^2\cdot 2^{-158}$.

\emph{What is open.} The assumption can be met: the argument whose proof is simply all the cells
is complete and knowledge sound, though its proof is as large as the matrix (an observation, not
formalized). The deployed argument is leanVM proving the relation inside the guest program. It cannot
meet the assumption in this model: its relation uses the real Keccak-f[1600], not $\pi$, so an
adversary computes a root with no query to $\pi$ and proves its opening honestly, and no
straight-line extractor finds a $\pi$-preimage of it. Code binding covers the arguments whose
relation is evaluated over $\pi$; for the deployed one it would need a model that ties the guest
program's Keccak to $\pi$ (Section~\ref{sec:notgive}, and the first item of the last section).

\section{The leanISA guest program (aggregation)}\label{sec:guest}

The guest program aggregates: from signatures, DA matrices and child proofs, it proves one
statement holding all their claims. This section defines what it proves and states what is proved
about it. Every definition named here is in \lean{Statement/LeanVmGuest.lean} unless marked.

\subsection{The objects}

\paragraph{The guest program} $G$ (\lean{LeanVM.Guest.guest}) is leanVM's compiled aggregation
program (\texttt{crates/rec\_aggregation/guests/lean\_ethereum.py} at \texttt{c445e3d5}),
transcribed as generated data (\lean{Statement/Data/}). Its bytecode $B$ has $2^{19}$ instruction
slots of $16$ words of $\mathbb{F} = \mathrm{GF}(2^{64})$, so $2^{23}$ words
(\lean{bytecode_length}, in \lean{lean_vm}). Below, $H$ is Keccak-256 and
$\mathbb{E} = \mathbb{F}[y]/(y^3+y+1)$, the field of $2^{192}$ elements. Its \emph{deployed seed}
\[
iv_G = \mathrm{IV}(B) = H(\mathit{TAG} \,\|\, \bar B), \qquad \mathit{TAG} = \texttt{leanvm-v1}
\]
(\lean{deployedSeed}, \lean{fiatShamirIvWith}, \texttt{verifier.py:684, 1464}), the hash of the
proof system's $9$-byte version tag and the bytecode's words as $8$-byte little-endian strings, is
the Fiat--Shamir initial value prover and verifier derive from the bytecode; a test run by the gate
(\texttt{lake exe node-vectors}) checks it against leanVM's \texttt{fs\_seed}.

\paragraph{A statement} $s$ (\lean{NodeStatement}) consists of
\begin{itemize}
\item a list $\Cx$ of leanXMSS claims and a list $\Cs$ of leanSPHINCS claims, each a pair $(pk, m)$
meaning ``$pk$ signed $m$'';
\item a \emph{DA digest} $\mathit{da}$ (\lean{Claims.da}): either $H(\rho \,\|\, \delta)$ of one
data-availability claim (\lean{DaClaim}, \lean{daClaimDigest}), a matrix root $\rho$
(\lean{daRoot}) and the digest $\delta$ (\lean{membershipVectorDigest}) of a vector $u$, or $32$
zero bytes without one (\lean{zeroDigest}; \texttt{da\_digest}, \texttt{aggregation.rs:550-562}).
The claim: some matrix of $1$ to $1024$ encoded blobs with root $\rho$ has every row orthogonal to
some $u$ with digest $\delta$ (\lean{Da.MatrixWitness.Opens}). The outer verifier holds $\rho$, derives
$\delta$ from $\rho$'s membership vector (Section~\ref{sec:schemes}), which makes the rows codewords
except with probability $14/2^{192}$, and checks $\mathit{da} = H(\rho \,\|\, \delta)$;
\item a \emph{deferred claim} $c = (z, v)$, $z \in \mathbb{E}^{23}$, $v \in \mathbb{E}$, meaning
that the multilinear extension $\widehat{B}$ of the bytecode, as a function on $\{0,1\}^{23}$,
satisfies $\widehat{B}(z) = v$ (\lean{DeferredClaim.Holds}).
\end{itemize}

\paragraph{The statement's digest} under a seed $\sigma$ (\lean{statementDigest},
\texttt{aggregation.rs:331-360}) is
\[
\mathrm{dig}_\sigma(s) = H\big(\sigma \,\|\, H(\enc(\Cx)) \,\|\, H(\enc(\Cs)) \,\|\, \mathit{da} \,\|\, \enc(c)\big),
\]
where $\enc(\Cx)$ concatenates, per claim, the key's root ($16$ bytes), its parameter ($16$) and
the message ($32$); $\enc(\Cs)$ likewise; and $\enc(c)$ is the $24$-byte encodings of
$z_1, \dots, z_{23}, v$,
padded with one zero element to an even number. The lists' lengths are bound by Keccak's padding;
the deferred claim's encoding carries no length, so deferred claims are told apart only at one
shape. The guest program's public input is this $32$-byte digest.

The seed $\sigma$ is the Fiat--Shamir IV with which $G$ verifies its children's proofs, which are
proofs of $G$ itself. $G$ cannot contain its own IV (a hash of its own bytecode), so it reads
$\sigma$ as input, and the digest binds it; the outer verifier computes the digest with
$\sigma = iv_G$.

\paragraph{A witness} $w$ (\lean{NodeWitness}) consists of
\begin{itemize}
\item a list of \emph{children}, each a statement $s_i$, a proof $p_i$, and a \emph{fresh} claim
$f_i = (z_i, v_i)$ on the bytecode, the evaluation the child's proof leaves to be checked;
\item for each signature claim of $s$, a \emph{holder}: a direct signature (leanXMSS or
leanSPHINCS) or the index of a child;
\item optionally, one LeanDA matrix the node checks directly, with the DA claim it opens;
\item the \emph{batching hints}: the round messages of one sumcheck, and a final value.
\end{itemize}

\subsection{What one node establishes}

$\mathrm{NodeRel}(B, \sigma, s, w)$ (\lean{NodeRelation}) is the conjunction of:
\begin{enumerate}
\item $s$'s deferred claim has the bytecode's shape: its point has $23$ coordinates
(\lean{DeferredClaim.HasBytecodeShape}).
\item Every child's deferred claim has the bytecode's shape, and the \emph{deferring verifier}
accepts its proof: $\mathrm{verifyDeferred}(B, \sigma, \mathrm{dig}_\sigma(s_i), f_i, p_i) =
\mathsf{true}$. The deferring verifier is the zkVM verifier except that, instead of evaluating
$\widehat{B}$ at the point it reaches, it outputs that evaluation as the claim $f_i$. It runs on
the real Keccak-256 and Keccak-f[1600], the ones the guest program computes.
\item Every signature claim of $s$ is held: a direct signature verifies under the claim's key and
message (Keccak-256 as the schemes' hash, at any leaf index); a child holder $i$ means the claim
occurs in the corresponding list of $s_i$.
\item The DA digest is accounted for: with a direct matrix, $\mathit{da} = H(\rho \,\|\, \delta)$ of the claim
it opens; without one, $\mathit{da}$ is zero or some child's DA digest (the guest program's
\texttt{da\_pick}, \texttt{lean\_ethereum.py:4207-4211}).
\item The deferred claim batches the children's. With no children, its point is all-zero (its
value is left to the outer verifier); otherwise $\mathrm{batch}\big((\mathrm{dig}_\sigma(s_i),
f_i, c_i)_i, \text{hints}\big) = c$ (\lean{batch}), with $c_i$ the deferred claim of $s_i$. The
batching is a Fiat--Shamir transcript on Keccak-f[1600], seeded from the hash of the label
\texttt{leanvm/recursion-aggregation/v1}: it absorbs the number of children, then each child's
digest and two claims; squeezes one challenge $\lambda \in \mathbb{E}$, whose powers batch the
claims (\lean{lambdaPowers}); and runs one degree-$2$ sumcheck reducing
$\sum_i (\lambda^{2i} v_i + \lambda^{2i+1} v'_i)$, with $v_i$
the value of $f_i$ and $v'_i$ that of $c_i$, to one claim on $\widehat{B}$ at a new point.
\end{enumerate}
\lean{NodeRelation} does \emph{not} say that a child's claims are true, only that the deferring verifier,
run on the real Keccak, accepts its proof. Whether that makes them true is the soundness of leanVM
with Keccak instantiated, which is not proved (Section~\ref{sec:notgive}).

\paragraph{The outer verifier} (\texttt{EthereumProof::verify}, \texttt{aggregation.rs:673-690})
additionally checks that the DA digest is that of its root with the digest of the root's
membership vector, that the deferred claim holds ($\widehat{B}(z) = v$), and the lists' form
(sorted, no repeats, fewer than $2^{16}$ signature claims). The soundness
statements below assume none of this: of the accepted statement they assume only that its deferred
claim has the bytecode's shape, which makes them stronger than statements about the outer
verifier.

\subsection{Soundness}

The outer verifier accepts a proof that $G$ ran on the public input $x = \mathrm{dig}_{iv_G}(s)$.
What does that establish about $s$? The answer has three parts, each an end theorem.
\begin{enumerate}
\item \textbf{The zkVM is knowledge sound} (\lean{KnowledgeSoundness},
Section~\ref{sec:zkvm}). Except with probability
\[
\varepsilon(Q) = (Q+2^{14})\cdot 2^{-128} + 5\,(Q+2^{14})^2\cdot 2^{-256},
\]
the extractor returns a valid execution $t$ of $G$ on $x$. This is the only step with an oracle
and a probability.
\item \textbf{Every execution of $G$ proves a statement} (\lean{GuestKnowledge}). From every
valid execution $t$ of $G$ on any public input $x$, the decoder $\mathrm{dec}$ reads a seed
$\sigma'$, a statement $s'$ and a witness $w'$ with
\[
\mathrm{dig}_{\sigma'}(s') = x \qquad\text{and}\qquad \mathrm{NodeRel}(B, \sigma', s', w').
\]
\item \textbf{The digest binds} (\lean{DigestBinds}). Let $(\sigma, s)$ and $(\sigma_0, s_0)$
be seeds and statements whose deferred claims have the bytecode's shape. If
$\mathrm{dig}_\sigma(s) = \mathrm{dig}_{\sigma_0}(s_0)$, then $(\sigma, s) = (\sigma_0, s_0)$, or
Keccak-256 collides on one of three pairs of inputs (\lean{collisionCandidates}): those of the
outer hash, of $H(\enc(\Cx))$ and of $H(\enc(\Cs))$. The DA digest enters the outer input as it
is, so it needs no pair of its own.
\end{enumerate}
Steps 2 and 3 are facts about $G$'s code and the digest: no oracle, no probability, and Keccak is
the deployed function $G$ computes.

\paragraph{Chained.} Steps 1 and 2 give \lean{aggregate_knowledge} (\lean{lean_vm}): except with
probability $\varepsilon(Q)$, an accepted proof of $G$ on any $x$ yields an execution whose
decoding $(\sigma', s', w')$ has $\mathrm{dig}_{\sigma'}(s') = x$ and
$\mathrm{NodeRel}(B, \sigma', s', w')$. Step 3 with $(\sigma_0, s_0) = (iv_G, s)$ then gives
\lean{AggregateSound}: except with probability $\varepsilon(Q)$, if the verifier accepts a proof
of $G$ on $\mathrm{dig}_{iv_G}(s)$, for any $s$ whose deferred claim has the bytecode's shape, the
extracted execution $t$ \emph{explains} $s$ (\lean{ExplainsStatement}):
\begin{itemize}
\item $\mathrm{dec}(t) = (iv_G, s, w')$ with $\mathrm{NodeRel}(B, iv_G, s, w')$, or
\item the decoded seed and statement, with $(iv_G, s)$, give a Keccak-256 collision among the
three pairs.
\end{itemize}
For one execution, step 3 alone is \lean{decode_binds} (\lean{lean_vm}).

Two Keccaks meet here. The verifier of step 1 runs on $\pi$: Keccak-256 is the sponge over $\pi$,
and Keccak-f is $\pi$. The digest, \lean{NodeRelation} (the children's proofs, the signatures, the
batching) and the collision use the deployed Keccak, which $G$ executes. The collision is therefore
not an event of small probability: it is a concrete pair of inputs, read off $t$ and $s$, on which
the deployed Keccak-256 collides.

\paragraph{Why a decoder.} \lean{GuestKnowledge} could have said only that some
$(\sigma', s', w')$ exists (\lean{guest_sound} proves that too). It says more: the witness is
\emph{in} the execution, and $\mathrm{dec}$ reads it. That is what makes the extracted execution
knowledge of a witness. But a decoder may use classical logic, and one that simply chose some
triple satisfying the conclusion would make \lean{GuestKnowledge} no stronger than the
existential. The statement is only as strong as its decoder, so a reviewer reads $\mathrm{dec}$
too.

The decoder is $\mathrm{dec} = $ \lean{LeanVM.Guest.End.decode}. It is defined in \lean{lean_vm},
not \lean{Statement/}, because it unfolds to $G$'s generated core code (about $484{,}000$ lines).
\lean{LeanVM.theorems_hold} proves \lean{GuestKnowledge} for it, and
\lean{LeanVM.corollaries_hold} proves \lean{AggregateSound} for it, from the end theorems. It reads
everything from the execution's memory:
\begin{itemize}
\item the seed and the statement from \texttt{main}'s frame;
\item each signature, matrix, child proof and batching hint from the frame of the function that
checked it: the first of the memory's $2^\kappa$ frames on which $G$'s own code for that function
holds for the call's arguments and results.
\end{itemize}
That frame test is $G$'s semantics, and it is the decoder's only search: it computes no hash and
runs no verifier. The test \lean{LeanVMTests.Guest.DecodeNoChoice} checks this by following every
definition the decoder unfolds to.

\paragraph{The batching on its own} (\lean{BatchingSound}). The batching of \lean{NodeRelation}'s last
condition is also proved sound as a verifier in its own right, in the RPM: Keccak-256 is the sponge
over $\pi$, Keccak-f is $\pi$, and it accepts when the sumcheck closes and the claim it outputs
holds of $\widehat{B}$ (\lean{runBatching} of \lean{boundedBatchingProgram}). Fix a bytecode $B$ of power-of-two
length and $M \le 2^{64}-1$. An adversary making $Q_A$ queries to $\pi$ and $\pi^{-1}$ outputs a
proof. The probability that the batching accepts it while its proof stream begins with fewer than
$M$ children of the bytecode's shape, one of whose claims (fresh or carried) is false
(\lean{SomeChildClaimFalse}), is at most
\[
\boxed{\ (Q_A + Q_V + 1)\cdot M\cdot 2^{-191} + 5\,(Q_A+Q_V)^2\cdot 2^{-256}\ }
\quad\text{for } Q_A \le 2^{128},\ Q_V \le 2^{96},
\]
where $Q_V$ (\lean{batchingQueryBound}) is the batching's own query bound for $M$ children (about
$2^{68}$ for $G$'s bytecode at $M = 2^{64}-1$; \lean{batch_sound_holds}). The factor $M$ comes from
the single challenge $\lambda$. Weighted by the powers of $\lambda$, the gaps of the fewer than $2M$
claims (claimed minus true value) form a polynomial in $\lambda$ of degree below $2M$, nonzero when
a claim is false. It vanishes at fewer than $2M$ values of $\lambda$, against at most $2$ values for
a sumcheck round's challenge. For $G$'s $16$ children, $M = 17$.

\subsection{What the soundness theorems do not give}\label{sec:notgive}

They speak about one node. To conclude from an accepted root that every claim held at the leaves,
one would need, at each internal node, the soundness of the children's proofs \emph{as the guest
program verifies them}: with the real Keccak-256 and Keccak-f, not with $\pi$. That is the
Fiat--Shamir heuristic for an instantiated hash, and it is not proved. Likewise \lean{BatchingSound} is
the batching outside the circuit, with Keccak-f $= \pi$; it does not compose with the in-circuit
statements, whose \texttt{KECCAK} instruction computes the deployed Keccak-f[1600], so soundness
across recursion levels does not follow from it.

\subsection{Completeness}

\paragraph{The hypotheses,} on a seed $\sigma$, a statement $s$ and a witness $w$ with
$\mathrm{NodeRel}(B, \sigma, s, w)$:
\begin{enumerate}
\item \lean{Limits}, the limits the guest program enforces on the node: a signature claim declared
$n$ times is held $n$ times, by direct signatures or by distinct entries of the children's lists;
at most $16$ children; each child's signature lists shorter than $2^{16}$; the coverage table (the
declared claims plus the duplicates the guest program records) shorter than $2^{16}$.
\item \lean{AcceptableChild} for every child, what the guest program asks of a child beyond
\lean{NodeRelation}:
  \begin{itemize}
  \item its proof announces a stack log of at least $22$ (the guest program has opening code for
  stack logs $22$ to $28$ only), and its proof stream has at most $8192$ field elements;
  \item at each query step of its deferring verifier, the Merkle openings agree wherever they share
  a node of the tree (an honest prover's openings, all from one tree, agree; disagreeing ones would
  have to collide Keccak-256);
  \item \emph{no zero challenge}: its proof is also accepted by the deferring verifier run through
  \lean{Prog.refusingZeroChallenge}, which fails on every zero challenge the verifier program samples
  (\lean{verifyDeferredRefusingZero}). The guest program divides by the fold and tail challenges and by the
  batching coefficient $\lambda$ at the stacked opening; refusing every zero challenge is slightly
  stronger, each challenge being zero with probability $2^{-192}$.
  \end{itemize}
  Honest children meet all of it, except with probability at most $2^{-126}$
  (\lean{GuestChildComplete} below).
\item The budget: a bound on the run's size, below.
\end{enumerate}

\paragraph{The budget} (\lean{lean_vm}). The node's \emph{work} (\lean{work}) consists of $\cnt{decl}$
declared claims (both signature lists), $\cnt{cover}$ coverage slots, $\cnt{xmss}$ and $\cnt{sphincs}$
signatures held directly, $\cnt{da} \le 1$ DA matrix held directly with $\cnt{rows}$ rows, padded to
$\cnt{pad}$ rows ($2^{\lceil\log_2 \ell\rceil}$ for a matrix of $\ell$ rows), $\cnt{child}$ children
with $\cnt{childclaims}$ claims in all,
and each child's announced stack log $\mu$ and rate exponent $r$. A run allocates at most
\begin{multline*}
W = 17789 + 35\cnt{decl} + 77\cnt{cover} + 6459\cnt{xmss} + 7328\cnt{sphincs} + 480944\cnt{da}
  + 212653\cnt{rows} \\ + 17672\cnt{pad} + 12941\cnt{child} + 30\cnt{childclaims}
  + \sum_{\text{children}} \beta(\mu, r)
\end{multline*}
words, where $\beta(\mu,r)$, the words of verifying one child, is:
\begin{center}\small
\begin{tabular}{l|rrrrrrr}
$\mu$ & 22 & 23 & 24 & 25 & 26 & 27 & 28\\\hline
$r=1$ & 562639 & 574166 & 596836 & 609461 & 620357 & 634168 & 656166\\
$r=2$ & 453149 & 458964 & 478503 & 485341 & 493109 & 500146 & 517970\\
$r=3$ & 411827 & 417258 & 434879 & 440555 & 446177 & 451897 & 469526\\
$r=4$ & 389193 & 393074 & 408904 & 413004 & 417160 & 423729 & 439992
\end{tabular}
\end{center}
The instructions of each opcode are bounded by linear forms of the same shape
(\lean{Count.deployedInstrs}; for \texttt{KECCAK}, $25 + \cnt{decl} + 398\cnt{xmss} + 390\cnt{sphincs} + 6047\cnt{da} +
2064\cnt{rows} + 129\cnt{pad} + 61\cnt{child} + \cnt{childclaims} + \sum \kappa(\mu, r)$ with
$\kappa$ between $4922$ and $9979$). The
budget's size is the committed size of Section~\ref{sec:zkvm} at these bounds: memory
$2^{\max(\lceil\log_2 W\rceil, 20)}$ and each table's height the least power of two at least its
bound (at least $2^9$ for \texttt{KECCAK}). These are proved bounds on the guest program's runs,
counted from its code, not measurements.

\paragraph{The theorems.} Each concludes that $G$ has a valid execution $t$ on public input
$\mathrm{dig}_\sigma(s)$ with $2^{22} \le \mathrm{size}(t) \le 2^{28}$ (committed size).
\begin{itemize}
\item \lean{GuestCompleteCovered}: under hypotheses 1 and 2 (the budget is discharged in the
proof), for every node in one of eight families (\lean{CoveredNode}), by its $\cnt{child}$ children,
$\cnt{xmss} + \cnt{sphincs}$ raw signatures, $\cnt{rows}$ blobs held directly (the rows of at most
one matrix), and claim counts $\cnt{decl}$, $\cnt{cover}$, $\cnt{childclaims}$ each at most the
family's cap:
\begin{center}\small
\begin{tabular}{lrrrr}
\toprule
Family & $\cnt{child}$ & $\cnt{xmss} + \cnt{sphincs}$ & $\cnt{rows}$ & claims\\\midrule
recursion & any & $0$ & $0$ & $2048$\\
leaf & $0$ & $\le 2048$ & $0$ & $2048$\\
& $\le 8$ & $0$ & $0$ & $2^{15}$\\
& $\le 4$ & $\le 1024$ & $0$ & $8192$\\
& $\le 4$ & $0$ & $\le 32$ & $8192$\\
& $\le 2$ & $\le 256$ & $\le 32$ & $8192$\\
& $\le 2$ & $0$ & $\le 32$ & $2^{14}$\\
& $0$ & $0$ & $\le 32$ & $2^{15}$\\
\bottomrule
\end{tabular}
\end{center}
Its first two families are the corollaries \lean{GuestCompleteRecursion} and
\lean{GuestCompleteLeaf}. \lean{Limits} caps a recursion node at $16$ children, each at
any $\mu \in [22,28]$ and $r \in [1,4]$.
\item \lean{guest_complete_numeric} (\lean{lean_vm}): under hypotheses 1 and 2, whenever the
budget's size is at most $2^{28}$. Evaluating the budget also places under $2^{28}$ a leaf of up to
$2550$ leanXMSS or $2252$ leanSPHINCS signatures (each declared once), and a node holding one DA
matrix directly with up to $57$ rows. On leanVM's own aggregations of $16$ and $32$ blobs the
budget's size is the measured committed size, $2^{26.486}$ and $2^{27.482}$ (at leanVM
\texttt{990b9352}, \texttt{93ebecc5} and \texttt{c445e3d5} alike).
\item \lean{aggregate_complete_of} (\lean{lean_vm}, in the RPM): the execution so obtained,
proved by the honest prover at $\rho \in [1,4]$ with $f \le 2^{64}$ nonces, is rejected with
probability at most $\delta(f)$ of Section~\ref{sec:zkvm}.
\end{itemize}

\paragraph{Honest proofs are children the guest program accepts} (\lean{GuestChildComplete}).
The hypotheses above concern the children. For the same named prover \lean{LeanVM.prover}:
\begin{enumerate}
\item in the RPM, with $2^{25}$ nonces per proof of work, the verifier that also refuses every zero
challenge (\lean{verifierRefusingZero}, through \lean{Prog.refusingZeroChallenge}) rejects its proof of any valid
execution of committed size at most $2^{28}$, at $1 \le \rho \le 4$, with probability at most
$2^{-126}$;
\item with the deployed Keccak (\lean{honestProofWithDeployedKeccak}), wherever \lean{verifierRefusingZero} accepts its proof of an
execution of $G$ of committed size at least $2^{22}$ on $\mathrm{dig}_{iv_G}(s)$, for $s$ of the
bytecode's shape, that proof is a child meeting \lean{AcceptableChild} for some fresh claim that is true of the
bytecode (\lean{BytecodeClaim.Holds}): the value the deferring verifier takes from the witness is
the bytecode's at the child's point.
\end{enumerate}
The guest program checks a child with the deployed Keccak, so the bad events (a proof of work with
no nonce within its fuel, a zero challenge) are bounded in the RPM for the same prover and check,
and their absence at the deployed Keccak gives the child's hypotheses deterministically. Completeness
thus holds level after level as far as the deployed Keccak behaves like $\pi$ for these events; that
step is the model's assumption, not a theorem. Not stated either: that honest batching hints exist,
that batching true claims gives a deferred claim that holds, and that the outer verifier accepts an
honest root (see the last section).

\section{What is not proved}

\begin{itemize}
\item The \texttt{KECCAK} instruction a program executes computes the real Keccak-f[1600], not
$\pi$: the RPM idealizes the proof system's hash, not the guest program's computation. It thus
excludes, by construction, programs that evaluate their own proof's Fiat--Shamir duplex
(Khovratovich, Rothblum and Soukhanov, 2025); the assumption is that Keccak-f[1600] behaves as a
random permutation even to programs that compute it. Statements
about the guest program's own execution use the deployed Keccak-256, under the assumption that no
one finds a Keccak-256 collision.
\item Soundness across recursion levels (Section~\ref{sec:notgive}); completeness across levels
holds only through \lean{GuestChildComplete}.
\item That the extractor runs in polynomial time (Section~\ref{sec:zkvm}); at deployed sizes, on
words with errors, it is infeasible, so the collision clause of \lean{AggregateSound} has no
concrete reduction cost.
\item Security against quantum adversaries (a quantum random-permutation model), for the
signature schemes and for leanVM's Fiat--Shamir: every statement is about classical adversaries.
\item LeanDA's code binding for an argument without access to $\pi$, and the link from the guest
program's DA check (\lean{Da.MatrixWitness.Opens}) and the outer verifier's membership-vector check to
LeanDA's code binding.
\item That a parent can always form its witness from honest children (the batching hints exist),
that batching true claims gives a true deferred claim, and that the deployed outer verifier
(\texttt{EthereumProof::verify}) accepts an honest root.
\item \lean{AcceptableChild}'s agreement of openings and absence of a zero challenge for children that are
not honest proofs.
\item The decoder \lean{LeanVM.Guest.End.decode}: a parameter of the statement, defined and read in
\lean{lean_vm}.
\item The Rust verifier, prover and guests. They are tested against the Lean specification: the
verifier differential fuzz; the prover differential (the Lean prover's proofs byte-identical to the
Rust prover's, accepted by both verifiers, every mutation rejected by both), which does not yet
include an execution of the guest program; the guest-diff check on real aggregations; and the node
vectors.
\item Witness generation: the Lean prover takes the execution (memory and hints) as input, checks
it, and computes the committed columns from it, but nothing in Lean runs a program on its private
inputs; in the prover differential the executions come from leanVM's Rust interpreter.
\item The Lean prover's speed: compiled, it takes about $18$ times as long as the multithreaded
Rust prover on the test corpus, and no complexity bound is proved.
\end{itemize}

\end{document}
